\section{Experiments}
\label{sec:experiments}

We design the experiments to answer five questions. \textbf{RQ1:} Which updates can be skipped, i.e., does the skip-harm controller preserve the full model's detections better than other selection policies at the same budget? \textbf{RQ2:} Why does it work? \textbf{RQ3:} Does it carry over to a different graph cell? \textbf{RQ4:} Which design choices matter? \textbf{RQ5:} Does selective updating reduce end-to-end latency? We also compare the full-update backbones.

\subsection{Dataset}
\label{sec:data}
We use SPRSound~\cite{zhang2022sprsound} with the data splits of the BioCAS 2024 event-detection challenge~\cite{zhang2024meta}: the official detection training and validation recordings for development, and the 2024 test set (1704 recordings) for testing. We use the later 2025 release (1309 recordings)~\cite{sprsound2025repo} only as an additional test set. All audio is 8\,kHz. An audit of the official split showed that the validation set shares 261 of its 266 patients with the training set. We therefore pooled the two, removed duplicates, and held out 15\% of the patients as a patient-disjoint validation set. The test sets are effectively patient-disjoint from training. Table~\ref{tab:data} summarizes the data.

\begin{table}[H]
\caption{Data splits. ``Adv.'' counts adventitious (abnormal) events.}
\label{tab:data}
\centering
\setlength{\tabcolsep}{3.5pt}
\begin{tabular}{lrrrrr}
\toprule
\textbf{Split} & \textbf{Recordings} & \textbf{Patients} & \textbf{Hours} & \textbf{Events} & \textbf{Adv.} \\
\midrule
Train & 2698 & 320 & 8.5 & 9952 & 2206 \\
Validation & 625 & 56 & 1.9 & 2260 & 661 \\
Test 2024 & 1704 & 324 & 5.9 & 7659 & 407 \\
Test 2025 & 1309 & 162 & 3.9 & 4706 & 2531 \\
\bottomrule
\end{tabular}
\end{table}

The test sets also follow a different annotation style from the training data (Table~\ref{tab:shift}). In training, events cover 45\% of a recording and are separated by gaps of about 0.4\,s. In the 2024 test set, events cover 76\% of the recording and almost touch each other. The 2025 set has the same style and a much higher share of adventitious events. Detectors trained on the training style are therefore penalized by strict onset/offset matching on the test sets, which keeps absolute event scores low for every model.

\begin{table}[H]
\caption{Annotation statistics differ between releases (durations in seconds; median [interquartile range]).}
\label{tab:shift}
\centering
\setlength{\tabcolsep}{4pt}
\resizebox{\columnwidth}{!}{%
\begin{tabular}{lccc}
\toprule
\textbf{Statistic} & \textbf{Train} & \textbf{Test 2024} & \textbf{Test 2025} \\
\midrule
Event coverage of recording & 45\% & 76\% & 68\% \\
Event duration & 1.30 [0.93, 1.76] & 2.06 [1.64, 2.51] & 1.92 [1.62, 2.31] \\
Gap between events & 0.39 [0.09, 1.01] & 0.04 [0.03, 0.06] & 0.08 [0.06, 0.12] \\
Adventitious share of events & 22\% & 5\% & 54\% \\
\bottomrule
\end{tabular}}
\end{table}

\subsection{Evaluation Metrics}
\label{sec:metrics}
\paragraph*{Event preservation} Our main metric asks whether selective updating keeps what the same model detects with full updates. We decode events from both the full-update and the selective posteriors, and compute the event-based F-score of the selective events against the full-update events ($F_{\mathrm{pres}}$). We also report $R^{\mathrm{adv}}_{\mathrm{pres}}$, the fraction of the full model's adventitious detections that the selective model also produces. These metrics do not depend on the annotation style. They measure agreement with the model's own detections, not correctness, so we always report them together with metrics against the annotations.

\paragraph*{Detection against the annotations} We report the event-based F-score ($F$) used by the challenge~\cite{zhang2024meta}, computed as in~\cite{mesaros2016metrics}, with DCASE-style collars~\cite{turpault2019sound}: 200\,ms at the onset, and the larger of 200\,ms and 20\% of the event duration at the offset. The official evaluator is not public, so these settings are ours. We also report the 1\,s segment-based F-score for adventitious classes (seg-adv-$F$), which is less sensitive to the annotation style, and the SPRSound harmonic score (HS) for classifying annotated segments~\cite{zhang2022grand}.

\paragraph*{Decoding} Posteriors are smoothed with a trailing moving average, so no future frames are used, and runs shorter than a minimum length are dropped. Both parameters are chosen on the validation set.

\subsection{Compared Methods}
\label{sec:policies}
\paragraph*{Selection policies} All policies run on the same trained weights, with the same greedy rule, budget accounting and refresh rule. We evaluate them at budgets of 12.5\%, 25\% and 50\% of nodes per step:
\emph{random}; \emph{periodic} (stalest regions first, i.e., round-robin); \emph{energy} (highest band-relative log-mel energy); \emph{acoustic change} (largest spectral change); \emph{change} (the learned head trained on representation change, i.e., AGN-Net's objective); two ablations of our target (\emph{hold-only} and \emph{label-free}); and our \emph{skip-harm} controller. As a reference, a \emph{hindsight oracle} computes the label-free utility by running the full update at every step. It is not available at inference time and is not budget-matched in computation. We also report the cheap path alone, with no graph updates.

\paragraph*{Full-update backbones} To check whether the graph state itself is useful, we compare models that update every step and share our front end, encoder, readout and training recipe: attention \emph{pooling} without memory; a causal dilated \emph{temporal convolutional network} (TCN; receptive field about 5\,s); a \emph{GRU}; a plain \emph{graph GRU} (neighbour-mean aggregation with a shared GRU cell); a \emph{frequency-convolution} cell (our cell with attention replaced by a convolution over nodes); and our acoustic \emph{AGN-Net}. All use the same fixed recipe without per-model tuning.

\subsection{Implementation Details}
All models are implemented in Python~3.11 with PyTorch~2.7.1 (CUDA~12.6) and trained on single GPUs (NVIDIA L4, L40S, A100 or H100). The event- and segment-based metrics use our own implementation, which is part of the released code. Every configuration is trained with three seeds. Tables report the mean and standard deviation over seeds. For paired comparisons we use a patient-level bootstrap with 2000 resamples on counts pooled over seeds. Bootstrap differences are computed on these pooled counts and can therefore differ slightly from differences of the seed means reported in the tables. Because this bootstrap treats the trained models as fixed, we also report whether a difference has the same sign in every seed. For non-inferiority of selective against full updating, we set a margin of $-0.01$ absolute event $F$ in advance. Selective updating is non-inferior if the lower bound of the 95\% interval of $F_{\mathrm{sel}}-F_{\mathrm{full}}$ is above this margin.

For latency we run batch-1 streaming on one pinned core of a central processing unit (CPU) server, with 60 test recordings per seed (about 20{,}000 chunks), and time each stage per chunk. We test two runtimes: a PyTorch implementation and a lightweight NumPy implementation of the same weights. Both perform true sparse updates (gather the selected nodes, compute on their subgraph, scatter back) and reproduce the batched model's outputs to within $3\times10^{-4}$. Configurations are interleaved per recording, and we report the median paired speedup over full updating on the same recording.

\paragraph*{Use of generative AI} Generative AI (Claude Opus~5.5, Anthropic) was used in this study to help write the experiment and analysis code, run the experiments on the computing cluster, analyse the results and draft the manuscript text. The author reviewed and verified all code, results and text (see Acknowledgments).
